\subsection*{Case 1: blue $-\omega$, $3-\omega$, $-1+\omega$, $2+\omega$}
\begin{figure}[h]\centering\includegraphics[width=\textwidth]{figures/ruleA_case1.png}\caption{Rule A, case 1: the row, the four forced points, and the boundary colouring of this case (ringed).}\end{figure}
\begin{center}\small\begin{longtable}{rlll}\toprule & point & colour & because\\\midrule\endhead
1 & $-1-\omega$ & red & $T(1)$ with $-2$, $-1$ (blue)\\
2 & $\omega$ & blue & $T(1)$ with $0$, $1$ (red)\\
3 & $1+\omega$ & blue & $T(1)$ with $1$, $2$ (red)\\
4 & $2-\omega$ & blue & $T(1)$ with $1$, $2$ (red)\\
5 & $7$ & red & $T(1,3,4)$ with $3$, $4$ (blue)\\
6 & $2\omega$ & red & $T(3)$ with $-\omega$, $3-\omega$ (blue)\\
7 & $3-2\omega$ & red & $T(1)$ with $2-\omega$, $3-\omega$ (blue)\\
8 & $-1+2\omega$ & red & $T(1)$ with $-1+\omega$, $\omega$ (blue)\\
9 & $1+2\omega$ & red & $T(1)$ with $1+\omega$, $2+\omega$ (blue)\\
10 & $-1+3\omega$ & blue & $T(1,3,4)$ with $-1-\omega$, $-1+2\omega$ (red)\\
11 & $4+2\omega$ & blue & $T(1,3,4)$ with $2\omega$, $1+2\omega$ (red)\\
\bottomrule\end{longtable}\end{center}
\paragraph{Subcase: $q(2+\omega)$ blue.}
\begin{center}\small\begin{longtable}{rlll}\toprule & point & colour & because\\\midrule\endhead
1 & $q(5-2\omega)$ & red & $T(1,3,4)$ with $-1+3\omega$, $q(2+\omega)$ (blue)\\
2 & $q(4-\omega)$ & blue & $T(1,3,4)$ with $q(5-2\omega)$, $7$ (red)\\
3 & $q(3\omega)$ & red & $T(1,3,4)$ with $-1+3\omega$, $q(4-\omega)$ (blue)\\
4 & $q(-\omega)$ & red & $T(1,3,4)$ with $2+\omega$, $q(4-\omega)$ (blue)\\
\bottomrule\end{longtable}\end{center}
\noindent\emph{Contradiction:} $q(3\omega)$, $0$, $q(-\omega)$ are all red and form $T(1,3,4)$.\par\medskip
\paragraph{Subcase: $q(2+\omega)$ red.}
\begin{center}\small\begin{longtable}{rlll}\toprule & point & colour & because\\\midrule\endhead
1 & $q(2)$ & blue & $T(1,3,4)$ with $q(2+\omega)$, $3-2\omega$ (red)\\
2 & $q(2-2\omega)$ & blue & $T(1,3,4)$ with $q(2+\omega)$, $3-2\omega$ (red)\\
3 & $q(3)$ & red & $T(1)$ with $q(2)$, $2+\omega$ (blue)\\
4 & $q(5-2\omega)$ & red & $T(1,3,4)$ with $q(2-2\omega)$, $4+2\omega$ (blue)\\
5 & $q(3\omega)$ & blue & $T(3)$ with $0$, $q(3)$ (red)\\
6 & $q(4-\omega)$ & blue & $T(1,3,4)$ with $q(5-2\omega)$, $7$ (red)\\
\bottomrule\end{longtable}\end{center}
\noindent\emph{Contradiction:} $q(3\omega)$, $-1+3\omega$, $q(4-\omega)$ are all blue and form $T(1,3,4)$.\par\medskip
\subsection*{Case 2: blue $3-\omega$, $-1+\omega$, $2+\omega$; red $-\omega$}
\begin{figure}[h]\centering\includegraphics[width=\textwidth]{figures/ruleA_case2.png}\caption{Rule A, case 2: the row, the four forced points, and the boundary colouring of this case (ringed).}\end{figure}
\begin{center}\small\begin{longtable}{rlll}\toprule & point & colour & because\\\midrule\endhead
1 & $-1-\omega$ & red & $T(1)$ with $-2$, $-1$ (blue)\\
2 & $\omega$ & blue & $T(1)$ with $0$, $1$ (red)\\
3 & $1+\omega$ & blue & $T(1)$ with $1$, $2$ (red)\\
4 & $2-\omega$ & blue & $T(1)$ with $1$, $2$ (red)\\
5 & $7$ & red & $T(1,3,4)$ with $3$, $4$ (blue)\\
6 & $3-2\omega$ & red & $T(1)$ with $2-\omega$, $3-\omega$ (blue)\\
7 & $-1+2\omega$ & red & $T(1)$ with $-1+\omega$, $\omega$ (blue)\\
8 & $1+2\omega$ & red & $T(1)$ with $1+\omega$, $2+\omega$ (blue)\\
9 & $-1+3\omega$ & blue & $T(1,3,4)$ with $-1-\omega$, $-1+2\omega$ (red)\\
10 & $2\omega$ & red & $T(1)$ with $\omega$, $1+\omega$ (blue)\\
11 & $4+2\omega$ & blue & $T(1,3,4)$ with $2\omega$, $1+2\omega$ (red)\\
\bottomrule\end{longtable}\end{center}
\paragraph{Subcase: $q(2+\omega)$ blue.}
\begin{center}\small\begin{longtable}{rlll}\toprule & point & colour & because\\\midrule\endhead
1 & $q(5-2\omega)$ & red & $T(1,3,4)$ with $-1+3\omega$, $q(2+\omega)$ (blue)\\
2 & $q(4-\omega)$ & blue & $T(1,3,4)$ with $q(5-2\omega)$, $7$ (red)\\
3 & $q(3\omega)$ & red & $T(1,3,4)$ with $-1+3\omega$, $q(4-\omega)$ (blue)\\
4 & $q(-\omega)$ & red & $T(1,3,4)$ with $2+\omega$, $q(4-\omega)$ (blue)\\
\bottomrule\end{longtable}\end{center}
\noindent\emph{Contradiction:} $q(3\omega)$, $0$, $q(-\omega)$ are all red and form $T(1,3,4)$.\par\medskip
\paragraph{Subcase: $q(2+\omega)$ red.}
\begin{center}\small\begin{longtable}{rlll}\toprule & point & colour & because\\\midrule\endhead
1 & $q(2)$ & blue & $T(1,3,4)$ with $q(2+\omega)$, $3-2\omega$ (red)\\
2 & $q(2-2\omega)$ & blue & $T(1,3,4)$ with $q(2+\omega)$, $3-2\omega$ (red)\\
3 & $q(3)$ & red & $T(1)$ with $q(2)$, $2+\omega$ (blue)\\
4 & $q(5-2\omega)$ & red & $T(1,3,4)$ with $q(2-2\omega)$, $4+2\omega$ (blue)\\
5 & $q(3\omega)$ & blue & $T(3)$ with $0$, $q(3)$ (red)\\
6 & $q(4-\omega)$ & blue & $T(1,3,4)$ with $q(5-2\omega)$, $7$ (red)\\
\bottomrule\end{longtable}\end{center}
\noindent\emph{Contradiction:} $q(3\omega)$, $-1+3\omega$, $q(4-\omega)$ are all blue and form $T(1,3,4)$.\par\medskip
\subsection*{Case 3: blue $3-\omega$, $2+\omega$; red $-\omega$, $-1+\omega$}
\begin{figure}[h]\centering\includegraphics[width=\textwidth]{figures/ruleA_case3.png}\caption{Rule A, case 3: the row, the four forced points, and the boundary colouring of this case (ringed).}\end{figure}
\begin{center}\small\begin{longtable}{rlll}\toprule & point & colour & because\\\midrule\endhead
1 & $3-4\omega$ & red & $T(4)$ with $-1$, $3$ (blue)\\
2 & $\omega$ & blue & $T(1)$ with $-1+\omega$, $0$ (red)\\
3 & $1+\omega$ & blue & $T(1)$ with $1$, $2$ (red)\\
4 & $7$ & red & $T(1,3,4)$ with $3$, $4$ (blue)\\
5 & $5-\omega$ & red & $T(1)$ with $4$, $5$ (blue)\\
6 & $4-5\omega$ & blue & $T(1,3,4)$ with $-\omega$, $3-4\omega$ (red)\\
7 & $1+2\omega$ & red & $T(1)$ with $1+\omega$, $2+\omega$ (blue)\\
8 & $2\omega$ & red & $T(1)$ with $\omega$, $1+\omega$ (blue)\\
9 & $4+2\omega$ & blue & $T(1,3,4)$ with $2\omega$, $1+2\omega$ (red)\\
\bottomrule\end{longtable}\end{center}
\paragraph{Subcase: $q(5-5\omega)$ blue.}
\begin{center}\small\begin{longtable}{rlll}\toprule & point & colour & because\\\midrule\endhead
1 & $q(1-\omega)$ & red & $T(4)$ with $4-5\omega$, $q(5-5\omega)$ (blue)\\
2 & $q(4-5\omega)$ & red & $T(1,3,4)$ with $4-5\omega$, $q(5-5\omega)$ (blue)\\
3 & $q(-\omega)$ & blue & $T(1)$ with $0$, $q(1-\omega)$ (red)\\
4 & $q(4-\omega)$ & blue & $T(4)$ with $q(4-5\omega)$, $7$ (red)\\
\bottomrule\end{longtable}\end{center}
\noindent\emph{Contradiction:} $q(-\omega)$, $2+\omega$, $q(4-\omega)$ are all blue and form $T(1,3,4)$.\par\medskip
\paragraph{Subcase: $q(5-5\omega)$ red.}
\begin{center}\small\begin{longtable}{rlll}\toprule & point & colour & because\\\midrule\endhead
1 & $q(5-\omega)$ & blue & $T(1,3,4)$ with $5-\omega$, $q(5-5\omega)$ (red)\\
2 & $q(5-2\omega)$ & blue & $T(3)$ with $q(5-5\omega)$, $7$ (red)\\
\bottomrule\end{longtable}\end{center}
\noindent\emph{Contradiction:} $q(5-\omega)$, $q(5-2\omega)$, $4+2\omega$ are all blue and form $T(1)$.\par\medskip
\subsection*{Case 4: blue $-\omega$, $2+\omega$; red $3-\omega$, $-1+\omega$}
\begin{figure}[h]\centering\includegraphics[width=\textwidth]{figures/ruleA_case4.png}\caption{Rule A, case 4: the row, the four forced points, and the boundary colouring of this case (ringed).}\end{figure}
\begin{center}\small\begin{longtable}{rlll}\toprule & point & colour & because\\\midrule\endhead
1 & $3-4\omega$ & red & $T(1,3,4)$ with $-1$, $-\omega$ (blue)\\
2 & $\omega$ & blue & $T(1)$ with $-1+\omega$, $0$ (red)\\
3 & $3-3\omega$ & blue & $T(1,3,4)$ with $-1+\omega$, $0$ (red)\\
4 & $1+\omega$ & blue & $T(1)$ with $1$, $2$ (red)\\
5 & $-1+3\omega$ & blue & $T(1,3,4)$ with $2$, $3-\omega$ (red)\\
6 & $6-4\omega$ & blue & $T(1,3,4)$ with $2$, $3-\omega$ (red)\\
7 & $7$ & red & $T(1,3,4)$ with $3$, $4$ (blue)\\
8 & $5-\omega$ & red & $T(1)$ with $4$, $5$ (blue)\\
9 & $1+2\omega$ & red & $T(1)$ with $1+\omega$, $2+\omega$ (blue)\\
10 & $2\omega$ & red & $T(1)$ with $\omega$, $1+\omega$ (blue)\\
11 & $4+2\omega$ & blue & $T(1,3,4)$ with $2\omega$, $1+2\omega$ (red)\\
\bottomrule\end{longtable}\end{center}
\paragraph{Subcase: $q(4-\omega)$ blue.}
\begin{center}\small\begin{longtable}{rlll}\toprule & point & colour & because\\\midrule\endhead
1 & $q(3\omega)$ & red & $T(1,3,4)$ with $-1+3\omega$, $q(4-\omega)$ (blue)\\
2 & $q(-\omega)$ & red & $T(1,3,4)$ with $2+\omega$, $q(4-\omega)$ (blue)\\
\bottomrule\end{longtable}\end{center}
\noindent\emph{Contradiction:} $q(3\omega)$, $0$, $q(-\omega)$ are all red and form $T(1,3,4)$.\par\medskip
\paragraph{Subcase: $q(4-\omega)$ red.}
\begin{center}\small\begin{longtable}{rlll}\toprule & point & colour & because\\\midrule\endhead
1 & $q(5-2\omega)$ & blue & $T(1,3,4)$ with $q(4-\omega)$, $7$ (red)\\
2 & $q(4-5\omega)$ & blue & $T(4)$ with $q(4-\omega)$, $7$ (red)\\
3 & $q(2+\omega)$ & red & $T(1,3,4)$ with $-1+3\omega$, $q(5-2\omega)$ (blue)\\
4 & $q(1-2\omega)$ & red & $T(4)$ with $-1+3\omega$, $q(5-2\omega)$ (blue)\\
5 & $q(2-2\omega)$ & red & $T(1,3,4)$ with $q(5-2\omega)$, $4+2\omega$ (blue)\\
6 & $q(5-\omega)$ & red & $T(1)$ with $q(5-2\omega)$, $4+2\omega$ (blue)\\
7 & $q(6-3\omega)$ & red & $T(1)$ with $q(5-2\omega)$, $4+2\omega$ (blue)\\
8 & $q(3-5\omega)$ & red & $T(1)$ with $q(4-5\omega)$, $6-4\omega$ (blue)\\
9 & $q(3)$ & blue & $T(1,3,4)$ with $q(2+\omega)$, $q(6-3\omega)$ (red)\\
10 & $q(1-\omega)$ & blue & $T(1)$ with $q(1-2\omega)$, $q(2-2\omega)$ (red)\\
11 & $q(5-5\omega)$ & blue & $T(1,3,4)$ with $q(5-\omega)$, $5-\omega$ (red)\\
12 & $q(3-4\omega)$ & red & $T(1,3,4)$ with $q(3)$, $2+\omega$ (blue)\\
13 & $4-5\omega$ & red & $T(4)$ with $q(1-\omega)$, $q(5-5\omega)$ (blue)\\
14 & $7-7\omega$ & blue & $T(1,3,4)$ with $q(3-4\omega)$, $q(3-5\omega)$ (red)\\
15 & $4-4\omega$ & blue & $T(1)$ with $3-4\omega$, $4-5\omega$ (red)\\
\bottomrule\end{longtable}\end{center}
\noindent\emph{Contradiction:} $3-3\omega$, $4-4\omega$, $7-7\omega$ are all blue and form $T(1,3,4)$.\par\medskip
